% Chapter 4 — The Machines. Author: KLEPPMANN. Budget 5 pp.
\section{The Machines}\label{ch:machines}

This chapter discharges constitution clauses C-5.1--C-5.8 (the three machines, the single-writer door, the cause-derived identifier, and the trust boundary).

The framework has three machines, one for each arrow of the typed picture of
Chapter~2 (\emph{The Picture: Map, Then Fold}): the Event Monitor evaluates
\texttt{watch}, the Events Executor evaluates \texttt{contract}, and the
Transaction Executor evaluates \texttt{apply}. All three are generic machinery.
No machine holds a line of product knowledge: what a dividend, a touch, or an
expiry \emph{means} lives in the smart contracts (Chapter~5); the machines only
notice, deliver, and admit. And only the last of the three ever writes. The
Event Monitor emits events; the Events Executor collects proposals; the
Transaction Executor alone changes the state of the ledger. Three roles, one
writer. In one picture --- machines in brackets, data on the arrows, the dotted
line the trust boundary:

{\small
\begin{verbatim}
 ................... outside the trust boundary ...................
 :                                                                :
 :      clock                          external events            :
 :        |                                  |                    :
 :        v                                  v                    :
 :   [Event Monitor] -- emits Event --> [Events Executor]         :
 :                                           |  fires             :
 :                                           v                    :
 :                                    [Smart Contract]            :
 :                                           |  proposes          :
 :                                           v                    :
 :                                      Transaction               :
 :...........................................|....................:
                                             v
               [Transaction Executor] -- admits --> Immutable Log
                    (the one writer)                     |
                                                         v
                                                    projections
\end{verbatim}
}

\subsection{The Event Monitor}

The Event Monitor owns time and attention. It holds every declared watch ---
the dates, observations, and conditions each unit's terms require watching
(Chapter~3) --- and evaluates them against the clock and the recorded
observations, emitting an event onto the record the moment a condition is met.
It holds no private list: the watches it evaluates are on the record, declared
by the units, so the Monitor's workload is itself a projection. And it emits
events, never transactions: nothing the Monitor does can place so much as one
move before the door.

The clock is the single input to the whole system that is not on the record,
and the Event Monitor is the single place it is consulted. The event the
Monitor emits \emph{is} recorded, so everything downstream of the emission is a
function of the record alone. Two consequences follow, and both are
load-bearing. A late firing produces the identical transaction: the contract
reads the recorded event and the recorded state, never the wall clock, so delay
costs timeliness and nothing else. And emission itself is verifiable: because
watches are declared data and the Monitor is deterministic, any party holding
the record and the clock recomputes which events should have been emitted, and
when --- the Monitor is auditable against its own inputs.

The relationship between the units and the Monitor is a recorded handshake.
Each declared watch is acknowledged by the Monitor, and the acknowledgement ---
an event like any other, flowing through the same pipeline --- is written into
the unit's state. A watch is thus visibly \emph{declared}, then \emph{armed},
then \emph{fired} or \emph{expired}, and each step is on the record. A declared
watch that is never armed is therefore a visible gap, never a silent one: it
stands in the unit's state as declared-but-unarmed, an overdue item any reader
can query. The full chain can be walked in either direction --- register a
unit, the conditions it declares, the acknowledgements, the events, the
transactions, the state updates --- which is the auditability commitment
(\S2) applied to attention itself.

The Monitor can affect timeliness and liveness; it cannot affect the
correctness of ledger state, because it only emits events, and an event changes
nothing until a contract's proposal passes the door.

\subsection{The Events Executor}

The Events Executor owns delivery. It does two things and nothing else. First,
it captures the events arriving from outside --- trades executed in the market,
corporate-action notices, fixings and other observations --- and \emph{proposes}
their recording: an observation crosses as a moveless observation-recording
transaction (Chapter~\ref{ch:marketdata}), which the Transaction Executor admits
or refuses like any other. Second, for each event, arrived or emitted, it fires the
responsible smart contract, handing it the event together with the current
projection of the ledger, and collects the proposed transaction. Many
contracts, one Events Executor: the contracts carry all the product knowledge;
the Events Executor carries none.

A trigger carries timing, not data: the firing says \emph{when}, and everything
the event means is read from the record. The Events Executor runs the
contracts; it never writes the ledger --- not even an observation --- and its
output is only proposals. A proposal is not a fact --- it becomes one only if the
one Transaction Executor admits it.

\subsection{The Transaction Executor}

The Transaction Executor owns the record. It is the gatekeeper of the ledger
door: the single component through which every transaction must pass, and the
only component permitted to change the ledger's state. On each proposal it
validates structure. Conservation is automatic --- a move is paired-leg by
construction, so the per-(unit, coordinate) sums cannot be disturbed
(Chapter~3); authorisation, idempotence, consistency of reference, and writer
discipline are enforced as checks at the door. If the proposal is admissible,
the Executor appends it to the immutable log; if not, it refuses, and the
refusal is a returned value, never a half-applied state. It creates nothing and
runs nothing. There is no other way to change the ledger's state.

Idempotence is the door's own defence, owed to no one upstream, and its key is the
transaction's own identifier. Start with the case the identifier must get right. One
corporate action --- a single cash dividend declared on stock ACME --- is one cause,
yet it touches every unit that references ACME, and each is adjusted by its own
transaction. The Monitor emits one event; the contract fires once per affected unit
and proposes one transaction per unit. These are \emph{different} transactions of the
\emph{same} cause, and every one of them must commit. An identifier that keyed on the
cause alone would collapse the whole set to one and silently drop every unit but the
first. So the cause alone cannot be the key.

\textbf{The derivation rule.} Every proposed transaction carries the identifier
\[
  \mathit{txid} \;=\; H(\mathit{causeEventId},\ \mathit{contractId},\ \mathit{unitId},\ \mathit{seq}),
\]
where $\mathit{causeEventId}$ is the recorded event that fired the contract,
$\mathit{contractId}$ and $\mathit{unitId}$ name the contract and the unit the
transaction acts on, $\mathit{seq}$ distinguishes the several legs one firing may
propose on a single unit, and $H$ is a fixed collision-resistant hash. The door keys
idempotence on $\mathit{txid}$: a proposal whose $\mathit{txid}$ already stands on the
log is committed zero further times; a proposal whose $\mathit{txid}$ is new is
admitted on its merits. Two properties make this rule exactly right, and each is
proved, not asserted.

\emph{Constant under retry.} Re-delivering the same firing recomputes the identical
tuple. All four components are read from the record: $\mathit{causeEventId}$ is the
recorded event, $\mathit{contractId}$ and $\mathit{unitId}$ are fixed by which
contract fired on which unit, and $\mathit{seq}$ is assigned by the contract as a pure
function of that recorded input (Chapter~\ref{ch:contracts}). A retry, a duplicate, or
a late re-firing therefore presents the same four components, so $H$ returns the same
$\mathit{txid}$; the door finds it already on the log and commits it once. Delay and
duplication cost timeliness and nothing else.

\emph{Injective over a cascade.} Fix one cause, so $\mathit{causeEventId}$ is held
constant across the whole cascade. The remaining three components
$(\mathit{contractId}, \mathit{unitId}, \mathit{seq})$ identify a leg uniquely: two
legs of the cascade differ in the unit they act on, or --- for two legs on one unit
--- in $\mathit{seq}$; no two distinct legs agree on all three. So the map
$(\mathit{contractId}, \mathit{unitId}, \mathit{seq}) \mapsto \mathit{txid}$ is
injective on the cascade's legs, \emph{provided} $H$ is injective on the tuple domain.
That is the standing assumption on $H$, stated rather than left implicit: over the
finite domain of well-formed cause tuples $H$ admits no collision --- exactly the
guarantee a collision-resistant hash is chosen to give. Under it, a cascade of $n$
legs carries $n$ distinct identifiers, the door sees $n$ new $\mathit{txid}$s, and all
$n$ commit. The two properties are jointly exact: a retry collapses because its tuple
is unchanged, and a genuine second leg commits because its tuple differs.

\begin{lstlisting}
txid :: EventId -> ContractId -> UnitId -> Seq -> Txid  -- txid = H(cause, contract, unit, seq)
-- retry:   same four arguments        => same Txid     => committed once (constant under retry)
-- cascade: one cause, distinct (contract, unit, seq)   => distinct Txids  => all n commit (injective)
\end{lstlisting}

The single-writer discipline is a correctness requirement, not an
implementation preference. The Transaction Executor presents to the log a
strictly monotonic sequence of committed transactions: every projection folds
the same total order, which is what makes two projections of one record unable
to disagree. Where two simultaneous events on the same instrument propose
transactions that do not commute and whose terms declare no precedence, the
Executor refuses rather than guesses: an order the record cannot justify is an
order nobody can recompute, and defaults fail closed (\S2). The refusal is
visible, and resolution --- a declared precedence, a corrected proposal ---
arrives through the same door.

\subsection{One event, three machines, one write}

This section runs OT-1 (Cast, \S\ref{sec:cast}) through the three machines and adds one
thing the pipeline walk of \S\ref{ch:picture} did not: what happens when the same emission
fires twice.

\textbf{Trigger.} OMEGA's official close of 79.00 is recorded. The declared
condition is met, and the Event Monitor emits the touch event onto the record.
No ledger state has changed: the Monitor emits events and cannot write state.

\textbf{Contract.} The Events Executor fires OT-1's smart contract, handing it
the touch event and the current projection. The contract reads the terms, the
holder (\mbox{W-ALPHA}), and the unit's state (\emph{live}), and returns one
transaction.

\textbf{Transaction.}

\begin{center}
\begin{tabular}{@{}rll@{}}
\toprule
 & Kind & Content \\
\midrule
1 & unit-state change & OT-1: UnitStatus \emph{live} $\to$ \emph{triggered} \\
2 & move & owned(USD) 1{,}000{,}000: W-BANK $\to$ W-ALPHA \\
\bottomrule
\end{tabular}
\end{center}

Where the unit stands encumbered under a collateral agreement, the payment leg
routes through that agreement's conditional obligation; Chapter~9 works this
same knock start to finish. The machines are identical in both tellings.

\textbf{The door.} The Transaction Executor checks the paired legs (the
per-(unit, coordinate) sums are undisturbed), the cause-derived identifier (the
touch event, not yet processed), and writer discipline (OT-1's state is written
by OT-1's contract), and admits. One append; the fold advances; the ledger now
shows the cash with \mbox{W-ALPHA} and the option \emph{triggered}.

\textbf{Duplicate emission.} Emission is at-least-once (Chapter~16), so the
Monitor may emit the touch again. The Events Executor fires the contract again;
the contract finds UnitStatus already \emph{triggered} and proposes nothing.
Had a duplicate proposal been produced anyway, its identifier recomputes to a
$\mathit{txid}$ already on the log and the door commits it zero further times. Two
independent defences, either one sufficient.

\emph{One event, three machines, one write --- and the recorded state is
exactly what makes a duplicate harmless.}

\subsection{The watch handshake}

\textbf{Trigger.} On 2026-05-04 the issuer of stock ACME announces a cash
dividend of 1.50 per share: record date 2026-05-15, payment date 2026-05-22.
The announcement is an external event; the Events Executor captures and records
it at the boundary.

\textbf{Contract.} The announcement fires ACME's smart contract once --- not to
pay, but to register the watches the dividend creates. Nothing is owed yet and
nothing moves; what the dividend has created is two future obligations of
attention, and the unit declares them exhaustively, exactly as it declared its
watches at registration.

\textbf{Transaction.} One moveless transaction:

\begin{center}
\begin{tabular}{@{}rll@{}}
\toprule
 & Kind & Content \\
\midrule
1 & unit-state change & ACME: record-date watch declared, 2026-05-15 \\
2 & unit-state change & ACME: payment-date watch declared, 2026-05-22 \\
\bottomrule
\end{tabular}
\end{center}

\textbf{The door.} A moveless transaction is well-formed: a transaction is a
bundle of four kinds of change and may carry no moves (Chapter~3). The
identifier derives from the recorded announcement; conservation is undisturbed
because nothing moved; the declaring writer is ACME's own contract.

\textbf{The handshake.} The Monitor acknowledges each declared watch; each
acknowledgement is an event like any other, flows through the same pipeline,
and lands as a unit-state update: both watches now visibly \emph{armed}. On
2026-05-15 the Monitor emits the record-date event and that watch is
\emph{fired}; on 2026-05-22 the payment-date watch fires likewise. What the two
firings do --- record each holder's entitlement movelessly, then pay against
the recorded entitlements --- is Chapter~5's episode, and the entitlement's
home is Chapter~\ref{ch:homes}'s; the machines' contribution is that neither
date could have passed silently. Had either watch remained unacknowledged
beyond its arming window, the unit's state would show a declared, unarmed
watch: an overdue item, never an absence.

\emph{A watch is visibly declared, then armed, then fired or expired; a date
that matters cannot pass silently, because attention itself is on the record.}

\subsection{The trust boundary}

The three roles sit at different levels of trust, and the dotted line of the
opening picture is exact: everything above it is untrusted. The Transaction
Executor is the trusted component --- the ledger's integrity rests on it and on
nothing else. The smart contracts, the Event Monitor, and the Events Executor
sit outside the boundary, untrusted in three different ways.

A contract can be economically wrong. Its proposal can be structurally
impeccable and still book the wrong economics, and the door cannot know that.
The defence is not prevention but recomputation: the events and the
transactions are both on the record, so verification re-runs the map on the
recorded event and compares (Chapter~15); a discrepancy is a defect of the
contract, not of the ledger, and it is repaired by a compensating transaction
through the same door.

The Event Monitor and the Events Executor cannot corrupt the ledger at all. A
late emission or a late firing produces the identical transaction, because
everything downstream reads the record and not the clock. A duplicate is
committed once, under the cause-derived identifier. A spurious firing finds
nothing due --- the one-touch's second touch above --- and proposes nothing, or
proposes something the door refuses. Every failure these two machines can
commit maps to delay or noise, never to wrong state.

What the untrusted machines do carry is liveness. Integrity owes them nothing;
timeliness owes them everything: watches must actually arm, events must
actually be emitted, contracts must actually be fired, obligations must
actually be discharged by their deadlines. Those duties are stated as minimum
requirements --- durable watches and timers, at-least-once emission,
acknowledged watches, triggers carrying timing not data, replayable
orchestration, no write privilege, obligation liveness --- in Chapter~16, and
each degrades, when unmet, to a visible overdue item rather than to silence.
